\paragraph{规范纯周期固定点的层计数、primitive 周期核与地址分形维数}
在定理 \ref{thm:xi-time-part9g-canonical-fixed-point-badic-periodic} 与推论 \ref{cor:xi-time-part9g-fixed-point-primitive-word-classification} 已把规范固定点的纯周期展开及其碰撞机制锁定之后，还可把同一地址系统的层计数、最小周期、全字母闭包与 prime-register 预算进一步压缩为如下五条闭合结论。

\begin{theorem}[规范纯周期固定点层的精确基数与前缀球分层]\label{thm:xi-time-part9gb-canonical-fixedpoint-count-prefix-stratification}
\leanverified{paper\_xi\_time\_part9gb\_canonical\_fixedpoint\_count\_prefix\_stratification}
沿用定理 \ref{thm:killo-godel-semidir-affine-2adic} 与定理
\ref{thm:xi-time-part9g-canonical-fixed-point-badic-periodic} 的记号。记
\[
T_v:=\ZZ_2v\subset T_2(E_{\min}),
\qquad
H_{\mathrm{can}}(k):=\{(r,k)\in H_{\mathrm{can}}\}
\qquad (k\ge 1).
\]
对每个
\[
h=(r,k)\in H_{\mathrm{can}}(k),
\]
记对应长度 \(k\) 的数字块为
\[
a(h):=\bigl(r(p_1),\dots,r(p_k)\bigr)\in\{0,\dots,M\}^k.
\]
则成立：
\begin{enumerate}
  \item \(\Psi(h)\) 在 \(T_v\) 上有且仅有一个固定点
  \[
  x_h=(1-B^k)^{-1}\operatorname{ev}_B(r)
  =
  \sum_{j\ge 0}\sum_{t=1}^{k}r(p_t)\,B^{jk+t-1}v.
  \]
  换言之，\(x_h\) 的 \(B\)-进展开恰为块 \(a(h)\) 的无限纯周期重复。
  \item 映射
  \[
  H_{\mathrm{can}}(k)\longrightarrow T_v,
  \qquad
  h\longmapsto x_h
  \]
  为单射，且
  \[
  \#\{x_h:\ h\in H_{\mathrm{can}}(k)\}=(M+1)^k.
  \]
  \item 对任意 \(0\le n\le k\) 及任意 \(h,h'\in H_{\mathrm{can}}(k)\)，有
  \[
  x_h-x_{h'}\in B^nT_v
  \quad\Longleftrightarrow\quad
  a(h)\ \text{与}\ a(h')\ \text{的前 }n\text{ 位相同}.
  \]
\end{enumerate}
\end{theorem}
\begin{proof}
第 \(1\) 条正是定理 \ref{thm:xi-time-part9g-canonical-fixed-point-badic-periodic} 的内容。

对第 \(2\) 条，由定理 \ref{thm:xi-time-part9g-affine-semigroup-unique-fixed-point}，在固定层级 \(k\) 上若 \(x_h=x_{h'}\)，则必有 \(h=h'\)，故 \(h\mapsto x_h\) 为单射。另一方面，\(H_{\mathrm{can}}(k)\) 的每个元素恰由 \(k\) 个彼此独立的 digits
\[
r(p_t)\in\{0,\dots,M\}
\qquad (1\le t\le k)
\]
唯一决定，因此
\[
|H_{\mathrm{can}}(k)|=(M+1)^k.
\]
于是固定点层的基数正为 \((M+1)^k\)。

再证第 \(3\) 条。对 \(h=(r,k)\)，记
\[
a(h)^\infty=(a_1,a_2,\dots)\in\{0,\dots,M\}^{\NN}
\]
为块 \(a(h)\) 的无限纯周期重复，则由第 \(1\) 条，
\[
x_h=\sum_{m\ge 1}a_mB^{m-1}v.
\]
对 \(h'\) 同理。应用定理 \ref{thm:xi-time-part9g-holographic-prefix-isometry-on-line} 于两条无限 digit 序列 \(a(h)^\infty\) 与 \(a(h')^\infty\)，得到
\[
x_h-x_{h'}\in B^nT_v
\quad\Longleftrightarrow\quad
a(h)^\infty,\ a(h')^\infty\ \text{的前 }n\text{ 位相同}.
\]
由于 \(0\le n\le k\)，上述前 \(n\) 位恰好就是块 \(a(h)\) 与 \(a(h')\) 的前 \(n\) 位，故结论成立。证毕。
\end{proof}

\begin{theorem}[规范纯周期固定点碰撞的 primitive 核与 Möbius 计数]\label{thm:xi-time-part9gb-fixedpoint-primitive-kernel-mobius}
\leanverified{paper\_xi\_time\_part9gb\_fixedpoint\_primitive\_kernel\_mobius}
沿用定理 \ref{thm:xi-time-part9gb-canonical-fixedpoint-count-prefix-stratification} 的记号。对任意长度 \(k\ge 1\) 的数字块
\[
a=(a_1,\dots,a_k)\in\{0,\dots,M\}^k,
\]
定义其对应纯周期点
\[
x(a):=\sum_{j\ge 0}\sum_{t=1}^{k}a_t\,B^{jk+t-1}v\in T_v.
\]
再记
\[
\mathcal F_k:=\{x_h:\ h\in H_{\mathrm{can}}(k)\}\subset T_v,
\]
以及最小周期层
\[
\Pi_k:=
\left\{
x\in \mathcal F_k:\ 
x\notin \mathcal F_d\ \text{对一切真因子 }d\mid k
\right\}.
\]
则：
\begin{enumerate}
  \item 对任意
  \[
  a\in\{0,\dots,M\}^k,
  \qquad
  b\in\{0,\dots,M\}^{\ell},
  \]
  有
  \[
  x(a)=x(b)
  \iff
  a^\infty=b^\infty
  \iff
  \exists\,u\ \text{为 primitive word},\ \exists\,r,s\ge 1:
  \ a=u^r,\ b=u^s.
  \]
  \item 对每个 \(k\ge 1\)，有不交并分解
  \[
  \mathcal F_k=\bigsqcup_{d\mid k}\Pi_d,
  \qquad
  |\mathcal F_k|=\sum_{d\mid k}|\Pi_d|.
  \]
  \item 若记
  \[
  P_M(k):=\#\{u\in\{0,\dots,M\}^k:\ u\ \text{为 primitive word}\},
  \]
  则
  \[
  |\Pi_k|=P_M(k)=\sum_{d\mid k}\mu(d)(M+1)^{k/d}.
  \]
  特别地，最小周期恰为 \(k\) 的规范纯周期固定点个数正是长度 \(k\) 的 primitive 字个数。
\end{enumerate}
\end{theorem}
\begin{proof}
第 \(1\) 条正是推论 \ref{cor:xi-time-part9g-fixed-point-primitive-word-classification} 的内容。

现证第 \(2\) 条。任取 \(x\in \mathcal F_k\)。由定理
\ref{thm:xi-time-part9gb-canonical-fixedpoint-count-prefix-stratification}，\(x\) 对应某个长度 \(k\) 的纯周期块 \(a\)。取 \(d\) 为无限字 \(a^\infty\) 的最小周期，则 \(d\mid k\)，并存在长度 \(d\) 的 primitive word \(u\) 使
\[
a=u^{k/d}.
\]
于是 \(x=x(u)\in \Pi_d\)。这表明
\[
\mathcal F_k\subseteq \bigsqcup_{d\mid k}\Pi_d.
\]
反向地，若 \(x\in \Pi_d\) 且 \(d\mid k\)，则 \(x=x(u)\) 对某个长度 \(d\) 的 primitive word \(u\) 成立；将其写成长度 \(k\) 的块 \(u^{k/d}\)，便得 \(x=x(u^{k/d})\in \mathcal F_k\)。故上述包含反向亦成立，且不同 \(d\) 的 \(\Pi_d\) 由最小周期定义两两不交，从而得到不交并分解及计数恒等式。

对第 \(3\) 条，定理 \ref{thm:xi-time-part9gb-canonical-fixedpoint-count-prefix-stratification} 已给出
\[
|\mathcal F_k|=(M+1)^k.
\]
与第 \(2\) 条合并得
\[
(M+1)^k=\sum_{d\mid k}|\Pi_d|.
\]
对该整除卷积施行 Möbius 反演，即得
\[
|\Pi_k|=\sum_{d\mid k}\mu(d)(M+1)^{k/d}.
\]
另一方面，第 \(1\) 条表明最小周期恰为 \(k\) 的点与长度 \(k\) 的 primitive words 一一对应，故 \(|\Pi_k|=P_M(k)\)。证毕。
\end{proof}

\begin{theorem}[全字母地址像的 \texorpdfstring{$2$}{2}-进自相似 Cantor 几何]\label{thm:xi-time-part9gb-full-alphabet-address-cantor-dimension}
沿用前述记号，定义全字母地址像
\[
\mathcal A_M
:=
\left\{
\sum_{t\ge 1}a_tB^{t-1}v:\ a_t\in\{0,\dots,M\}
\right\}
\subset T_v\subset T_2(E_{\min}),
\]
以及规范固定点全集
\[
\operatorname{Fix}_{\mathrm{can}}
:=
\bigcup_{k\ge 1}\mathcal F_k.
\]
则成立：
\begin{enumerate}
  \item \(\operatorname{Fix}_{\mathrm{can}}\) 在 \(\mathcal A_M\) 中稠密，亦即
  \[
  \overline{\operatorname{Fix}_{\mathrm{can}}}=\mathcal A_M.
  \]
  \item \(\mathcal A_M\) 是由 \((M+1)\) 个仿射收缩
  \[
  f_d(x):=dv+Bx
  \qquad (d\in\{0,\dots,M\})
  \]
  生成的严格自相似紧集，并满足不交分解
  \[
  \mathcal A_M=\bigsqcup_{d=0}^{M}f_d(\mathcal A_M).
  \]
  \item 若 \(M\ge 1\)，则 \(\mathcal A_M\) 是 \(T_v\) 中的紧、全不连通且无孤立点的 Cantor 型子集；并且在 \(T_2(E_{\min})\) 中的 Hausdorff 维数与 Minkowski 维数一致，满足
  \[
  \dim_{\mathrm H}(\mathcal A_M)
  =
  \dim_{\mathrm M}(\mathcal A_M)
  =
  \frac{\log(M+1)}{L\log 2}.
  \]
  此外，\(\mathcal A_M\) 在 \(T_2(E_{\min})\) 的 Haar 测度下为零测集。
\end{enumerate}
\end{theorem}
\begin{proof}
先证第 \(1\) 条。任取
\[
x=\sum_{t\ge 1}a_tB^{t-1}v\in \mathcal A_M.
\]
对每个 \(k\ge 1\)，把前 \(k\) 个 digits 周期化，定义
\[
x^{(k)}
:=
\sum_{j\ge 0}\sum_{t=1}^{k}a_t\,B^{jk+t-1}v.
\]
由定理 \ref{thm:xi-time-part9gb-canonical-fixedpoint-count-prefix-stratification}，\(x^{(k)}\in\mathcal F_k\subset \operatorname{Fix}_{\mathrm{can}}\)。同时 \(x^{(k)}\) 与 \(x\) 的前 \(k\) 个 \(B\)-进 digits 完全一致，故
\[
x^{(k)}-x\in B^kT_v.
\]
由于 \(k\to\infty\) 时 \(B^kT_v\) 在 \(2\)-进拓扑下收缩到 \(\{0\}\)，遂得 \(x^{(k)}\to x\)。这表明 \(\mathcal A_M\subseteq \overline{\operatorname{Fix}_{\mathrm{can}}}\)。反向包含显然，因为每个规范固定点本身都属于 \(\mathcal A_M\)。故闭包等式成立。

对第 \(2\) 与第 \(3\) 条，只需把定理 \ref{thm:xi-time-part62d-hologram-image-cantor-homeomorphism} 与定理 \ref{thm:xi-time-part62d-hologram-image-zero-haar-dimension} 应用于有限事件集
\[
E=\{0,\dots,M\},
\qquad
\mathrm{code}=\mathrm{id}_E.
\]
此时对应像集 \(\Lambda_{E,v}\) 恰为 \(\mathcal A_M\)。于是严格自相似分解、紧性、全不连通性、\(M\ge 1\) 时的无孤立点、维数公式以及 Haar 零测性同时成立。证毕。
\end{proof}

\begin{theorem}[Zeckendorf--Gödel 地址像的 Fibonacci 分形维数]\label{thm:xi-time-part9gb-zg-address-fractal-dimension}
取二值字母表 \(\{0,1\}\)，并定义 Zeckendorf--Gödel 地址像
\[
\mathcal A_{\mathrm{ZG}}
:=
\left\{
\sum_{t\ge 1}z_tB^{t-1}v:\ 
z_t\in\{0,1\},\ z_tz_{t+1}=0
\right\}
\subset T_v.
\]
则 \(\mathcal A_{\mathrm{ZG}}\) 正是禁止相邻 \(1\) 子移位在地址线上的 \(2\)-进自仿射吸引子，并满足：
\begin{enumerate}
  \item \(\mathcal A_{\mathrm{ZG}}\) 满足严格自仿射方程
  \[
  \mathcal A_{\mathrm{ZG}}
  =
  B\mathcal A_{\mathrm{ZG}}
  \ \sqcup\
  \bigl(v+B^2\mathcal A_{\mathrm{ZG}}\bigr).
  \]
  \item 在尺度 \(B^{-n}\) 上，\(\mathcal A_{\mathrm{ZG}}\) 恰分解为 \(F_{n+2}\) 个两两不交的 cylinder，其中 \(F_n\) 为 Fibonacci 数列。
  \item \(\mathcal A_{\mathrm{ZG}}\) 在 \(T_2(E_{\min})\) 中的 Hausdorff 维数与 Minkowski 维数一致，且
  \[
  \dim_{\mathrm H}(\mathcal A_{\mathrm{ZG}})
  =
  \dim_{\mathrm M}(\mathcal A_{\mathrm{ZG}})
  =
  \frac{\log\varphi}{L\log 2}.
  \]
\end{enumerate}
\end{theorem}
\begin{proof}
第 \(1\) 条正是定理 \ref{thm:xi-zg-address-binary-self-affine} 的自仿射分解。

对第 \(2\) 与第 \(3\) 条，定理 \ref{thm:xi-zg-address-minkowski-dimension-content} 已证明：在地址线 \(T_v\) 上，\(\mathcal A_{\mathrm{ZG}}\) 于尺度 \(B^{-n}\) 的最少球覆盖数恰为
\[
N_{B,v}(n)=F_{n+2},
\]
并且
\[
\dim_{\mathrm H}(\mathcal A_{\mathrm{ZG}})
=
\dim_{\mathrm M}(\mathcal A_{\mathrm{ZG}})
=
\frac{\log\varphi}{\log B}
=
\frac{\log\varphi}{L\log 2}.
\]
最后，由定理 \ref{thm:xi-time-part62d-hologram-image-zero-haar-dimension} 证明中关于 rank-one 地址线嵌入的双 Lipschitz 论证可知，\(\mathcal A_{\mathrm{ZG}}\) 在环境空间 \(T_2(E_{\min})\) 中的 Hausdorff 与 Minkowski 维数，等于其在地址线 \(T_v\) 中的相应维数。因此上述数值在 \(T_2(E_{\min})\) 中保持不变，结论成立。证毕。
\end{proof}

\begin{corollary}[prime-register 最坏层成本的地址维数补码解释]\label{cor:xi-time-part9gb-prime-register-dimension-complement}
若满足定理 \ref{thm:xi-layered-prime-register-sharp-natural-extension} 的最坏情形假设，即每一层都存在真实分支，则分层 prime-register 外置的最小层成本恰为每层一枚有效 prime slice。

把每一层 \(B=2^L\) 的地址槽位视为总宽度 \(L\) 比特，则：
\begin{enumerate}
  \item 对全字母地址像 \(\mathcal A_M\)，长度 \(k\) 的 cylinder 数为 \((M+1)^k\)，故层信息率为
  \[
  \lim_{k\to\infty}\frac{1}{k}\log_2 (M+1)^k=\log_2(M+1).
  \]
  由定理 \ref{thm:xi-time-part9gb-full-alphabet-address-cantor-dimension}，
  \[
  L\,\dim_{\mathrm H}(\mathcal A_M)=\log_2(M+1),
  \]
  从而其地址维数补码恰为
  \[
  L\bigl(1-\dim_{\mathrm H}(\mathcal A_M)\bigr)
  =
  L-\log_2(M+1).
  \]
  \item 对 Zeckendorf--Gödel 地址像 \(\mathcal A_{\mathrm{ZG}}\)，长度 \(k\) 的允许词数为 \(F_{k+2}\)，故层信息率为
  \[
  \lim_{k\to\infty}\frac{1}{k}\log_2 F_{k+2}=\log_2\varphi.
  \]
  由定理 \ref{thm:xi-time-part9gb-zg-address-fractal-dimension}，
  \[
  L\,\dim_{\mathrm H}(\mathcal A_{\mathrm{ZG}})=\log_2\varphi,
  \]
  因而其地址维数补码恰为
  \[
  L\bigl(1-\dim_{\mathrm H}(\mathcal A_{\mathrm{ZG}})\bigr)
  =
  L-\log_2\varphi.
  \]
\end{enumerate}
因此，在最坏情形下必须逐层引入的 prime-register 外置，可在地址几何上解释为：它所补偿的正是每个 \(B\)-进层片中未被内部符号分裂吸收的维数缺口。
\end{corollary}
\begin{proof}
定理 \ref{thm:xi-layered-prime-register-sharp-natural-extension} 已给出最坏情形下一层一枚有效 prime slice 的必要且可达下界。

对 \(\mathcal A_M\)，长度 \(k\) 的可用地址前缀恰为全部 \(k\)-位 blocks，故其每层信息率为 \(\log_2(M+1)\) 比特；再由定理 \ref{thm:xi-time-part9gb-full-alphabet-address-cantor-dimension} 的维数公式，
\[
\dim_{\mathrm H}(\mathcal A_M)=\frac{\log(M+1)}{L\log 2},
\]
两式等价于
\[
L\,\dim_{\mathrm H}(\mathcal A_M)=\log_2(M+1).
\]
于是补码公式立即成立。

对 \(\mathcal A_{\mathrm{ZG}}\)，定理 \ref{thm:xi-time-part9gb-zg-address-fractal-dimension} 的第 \(2\) 条给出长度 \(k\) 的允许词数为 \(F_{k+2}\)，由 Fibonacci 增长率
\[
F_{k+2}\sim \frac{\varphi^{k+2}}{\sqrt5}
\]
可得其每层信息率为 \(\log_2\varphi\)。再把同一定理中的维数公式
\[
\dim_{\mathrm H}(\mathcal A_{\mathrm{ZG}})=\frac{\log\varphi}{L\log 2}
\]
改写为
\[
L\,\dim_{\mathrm H}(\mathcal A_{\mathrm{ZG}})=\log_2\varphi,
\]
便得到第二条补码公式。证毕。
\end{proof}

\endinput
